\worksheettitle{Выбираем по разрядам}{\modulemark{M06-R} Коррекционная карточка}
\studentnameblock

\begin{notebox}[Задача]
  Из цифр \(0,\ 3,\ 4,\ 9\) составьте максимум и минимум. Каждую цифру
  используйте ровно один раз.
\end{notebox}

\begin{practicebox}[1. Максимум: слева направо]
  Тысячи: наибольшая цифра \answerblank[14mm].
  Сотни: наибольшая из оставшихся \answerblank[14mm].

  Десятки: \answerblank[14mm]. Единицы: \answerblank[14mm].
  Максимум: \answerblank[40mm].
\end{practicebox}

\begin{practicebox}[2. Минимум: ноль не первый]
  Тысячи: наименьшая ненулевая цифра \answerblank[14mm].
  Сотни: теперь можно поставить \answerblank[14mm].

  Остальные цифры по возрастанию: \answerblank[14mm], \answerblank[14mm].
  Минимум: \answerblank[40mm].
\end{practicebox}

\begin{reflectionbox}[3. Сравните по разрядам]
  Почему \(3049<3409\)? Тысячи одинаковы. Первый различающийся разряд \textemdash{}
  \answerblank[23mm]; в нём \answerblank[14mm] меньше \answerblank[14mm].
\end{reflectionbox}

\begin{warningbox}[Проверка условий]
  В каждом ответе должно быть столько цифр, сколько дано в наборе. Ни одна
  цифра не теряется и не используется лишний раз. Запись числа не начинается
  с нуля.
\end{warningbox}

\newpage
\section{Повторяем самостоятельно}

\begin{practicebox}[4. Несколько нулей]
  Цифры \(0,\ 0,\ 5,\ 7\).
  \begin{enumerate}[label=\textbf{\arabic*.}]
    \item Максимум: \answerblank[35mm]. Почему нули стоят справа?
      \writinglines{1}
    \item Минимум: \answerblank[35mm]. Почему первая цифра равна 5?
      \writinglines{1}
  \end{enumerate}
\end{practicebox}

\begin{practicebox}[5. Повторяющаяся цифра]
  Цифры \(0,\ 2,\ 2,\ 6\). Максимум: \answerblank[34mm]. Минимум:
  \answerblank[34mm]. Подчеркните в каждом числе обе цифры 2 и проверьте, что
  ни одна цифра не потеряна.
\end{practicebox}

\begin{warningbox}[6. Найдите ошибку]
  Для цифр \(0,\ 2,\ 6,\ 8\) минимум записали \(2608\). Правильный ответ:
  \answerblank[34mm]. Первый различающийся разряд в числах \(2068\) и
  \(2608\): \answerblank[24mm]. Объяснение: \answerblank[62mm].
\end{warningbox}

\begin{checkpointbox}[Повторная проверка]
  Новый набор цифр: \(0,\ 3,\ 3,\ 7\).
  \begin{enumerate}[label=\textbf{\arabic*.}]
    \item Максимум: \answerblank[34mm].
    \item Минимум: \answerblank[34mm].
    \item Объясните положение нуля в минимуме и проверьте, что обе цифры \(3\)
      использованы.
      \writinglines{2}
  \end{enumerate}
\end{checkpointbox}
